\section{SRAM DCIM：原生单tile数字求值与普通写入}

\subsection{状态、原生结构与读码生命周期}

本例采用 D6CIM 型二进制6T SRAM和精确整数数字乘加。\sid{SDCIM-01} 的 $128\times128$ 物理bit阵列对应 $K=128,N=16$ 的INT8权重矩阵，每八列共用一个HCA混合压缩加法树和BFA位优先移位累加器，共16条输出通道。一个完整MAC时钟处理一位输入、八位权重、16个输入项；八个输入bit之后更换行组。原生完整向量因此需要64个时钟，输出16个23-bit结果（PDF pp.1--3，Figs.1、4、5）。

参考配置直接采用一个原生tile，16384个数据cell表示2048 Byte有效状态；输入128 Byte、寄存1024 bit，输出寄存 $16\times23=368$ bit。每权重八个cell，内部互补节点不增加payload；16条已安装HCA/BFA均活动。一个控制器和更新域服务全部状态。所选16项归约保留原始时钟证据，不将其改为32项并沿用时序。

静态双字线连接使同组权重跨八个输入位保持，输入bit按位优先顺序计算。每向量有八次行组选择、64个MAC槽，不需要新增4096-bit操作数缓存。权重保持在原静态局部连接中，BFA累加器独立保存部分和，完成本组八个位后才切到下一行组。公共接口以 \texttt{static\_connection} 表示已有保持，完整MAC槽已含读服务，所以额外 \texttt{read\_ns}=0；该零值表示阶段覆盖。

\subsection{原文时间与完整周期覆盖}

D6CIM Fig.10(a)（PDF p.4）在0.9 V标示233 MHz，对应4.292 ns，采用向上取整4.3 ns完整MAC槽。参考和长槽为5/10 ns，与共同数字预算相容；短场景的边界捕获/提交各2 ns，MAC仍使用4.3 ns。原文1.1 V正文360 MHz与图363 MHz保留原值，不用于外推0.9 V。三个场景固定一个tile、16条通道和16项归约。

一个MAC槽覆盖静态SRAM连接、NOR乘法、HCA归约及BFA符号、移位、累加，不另加普通read access或ADC。补码符号bit取负位权；128项INT8累加的保守容器为23 bit，符合原输出宽度。

普通写周期独立于计算时钟。\sid{SDCIM-01} Fig.1给出128-bit同步memory接口，但未单独报告写完成时间。\sid{CMOS-07} 的28 nm eFlash平台48-bit、1k-word、2RW 8T SRAM在Fig.2展示同沿命令/地址/数据捕获及连续写，Table II典型455 MHz对应2.1978 ns完整时钟；1.54 ns是read access，不能作为写周期。典型供电1.05 V与本计算锚点不同。\sid{CMOS-03} 的28 nm 6T、TT/1 V、128-cell/15 fF瞬态模型采用约1 ns完整周期，以周期末节点达到目标并可连续访问为成功终点。

采用完整同步写槽 $T_M=2.2/5/10$ ns；2.2 ns由455 MHz向上取整，5/10 ns为正常完整周期预算。固定128对真实BL驱动、行译码、数据锁存及相应供电，不因场景改变资源。48-bit 8T/eFlash到128-bit 6T的移植仍是主要工程资格条件，不能称为D6CIM同芯片写实测。

\subsection{完整向量与resident装载}

本地整向量捕获一拍，八个行组各执行八输入bit，最后全部结果提交一拍；没有跨向量流水。显式形状代入公共接口为
\begin{equation}
N_D=8\times1\times\lceil128/16\rceil\times\lceil16/16\rceil=64,
\qquad\Delta_S=64T_C+2T_D.
\label{02_sram_dcim:eq:dcim-local-stream}
\end{equation}
内部位展开只计占用，不重复增加128 Byte输入。典型 $\Delta_S=64\times5+10=330$ ns。

每笔对齐行写128个cell，完整完成16个INT8权重。同步端口从捕获命令/地址/数据到节点翻转、关字线及恢复为下一次可计算状态，完整槽已覆盖本地控制；不再加同一事务的两拍准备/完成。二进制SRAM不要求擦除、迭代program/verify、破坏读恢复或定时刷新。因此
\begin{equation}
B_R=16\ \mathrm{Byte},\quad\Delta_R=T_M,\quad
T_R=128T_M,\quad\tau=16/T_M.
\label{02_sram_dcim:eq:dcim-local-resident}
\end{equation}
完整矩阵由128个16 Byte事务覆盖一次，典型 $T_R=640$ ns。小于满接口的零散写应按实际有效payload与所占完整周期计算，不领取16 Byte满口能力。掉电后重载属于新resident需求，寿命不进入瞬时吞吐。

\subsection{成对结果、组织对照与两表接口}

\begin{equation}
\rho=\frac{128}{64T_C+2T_D},\qquad
\mathrm{RI}^*=\frac{8T_M}{64T_C+2T_D},\qquad
U^*=\frac{128T_M}{64T_C+2T_D}=16\mathrm{RI}^*.
\label{02_sram_dcim:eq:dcim-local-ridge}
\end{equation}
时间以ns代入得到十进制GB/s。完整装载与完整向量的匹配关系已由公共映射接口保存。

% Generated by scripts/check_sram_dcim.py --emit from data/inputs.json.
\begin{table}[htbp]\centering\small
\begin{tabular}{@{}lrrrrrrr@{}}\toprule
情景 & $T_D$ & $T_C$ & $\Delta_S$ & $\Delta_R$ & $\rho$ & $\tau$ & $\mathrm{RI}^{*}$\\\midrule
短时隙 & 2 & 4.3 & 279.2 & 2.2 & 458.5 & 7273 & 0.06304\\
参考 & 5 & 5 & 330 & 5 & 387.9 & 3200 & 0.1212\\
长时隙 & 10 & 10 & 660 & 10 & 193.9 & 1600 & 0.1212\\
\bottomrule\end{tabular}
\caption{D6CIM型16项结构的成对情景。所有时间单位为ns，吞吐为十进制MB/s；$\Delta_R=T_M$。时间列保留复算值，非实测有效位数。}\label{02_sram_dcim:tab:dcim-results}\end{table}


典型 $\rho\approx0.39$ GB/s、$\tau=3.2$ GB/s、$\mathrm{RI}^*\approx0.12$，$U^*=640/330\approx1.94$ 个向量。Streaming主要受64个完整MAC槽约束，边界仅10/330约3\%；写侧由一个128-bit完整同步周期决定。参考和长场景计算、写周期同比变化，RI*相同；短场景的MAC受原4.3 ns证据约束，而写为2.2 ns，故两路比例不同。该比例不直接给出能效或介质总体排序。

八个原生tile可表示 $K=128,N=128$ 的较大逻辑矩阵。如果安装128条HCA/BFA却只令16条同时活动，并由单写域串行服务，完整向量需512个MAC槽、整矩阵1024次写。该明确容量组织典型 $\Delta_S=2570$ ns、$T_R=5120$ ns、$U^*\approx1.99$，与原生单tile不同。输入共享使边界只付一次，因此扩展不完全等比例。表中的对照保留16项归约，不作未经验证的32项同钟假设。

% Generated by scripts/check_sram_dcim.py --emit.
\begin{table}[htbp]\centering\small
\begin{tabular}{@{}lrrrrr@{}}\toprule
参考时隙条件 & $\Delta_S$ (ns) & $\Delta_R$ (ns) & $\rho$ (MB/s) & $\tau$ (MB/s) & $\mathrm{RI}^{*}$\\\midrule
主情景：16项，完整写周期 & 330 & 5 & 387.9 & 3200 & 0.1212\\
八个原生tile串行 & 2570 & 5 & 49.81 & 3200 & 0.01556\\
另有两拍外部写握手 & 330 & 15 & 387.9 & 1067 & 0.3636\\
\bottomrule\end{tabular}
\caption{结构及边界条件的独立影响。八个原生tile保留16项归约与同一写域；最后一行只适用于确有额外握手的实现。}\label{02_sram_dcim:tab:dcim-sensitivity}\end{table}


额外任务层握手仅在确实存在时另计。若同步写之外另有准备/提交各一拍，典型局部写为15 ns，tau成为原来的三分之一；它是边界条件对照，不能与已完整覆盖的主服务相加两次。

计算和写入共享状态与控制，两路能力分别服务时取得，未假设可同时达到峰值。将该配置映射到Table II其他N、输入共享或共享更新组织时，应重新给出完整 $T_R,\Delta_S$，再求 $U^*$。

\subsection{参数证据与适用范围}

原生归约、静态保持和64拍完成都有D6CIM结构依据。剩余关键条件是共同外围下的完整MAC预算，以及128-bit 6T普通写周期的跨实现移植；输入、原值/采用值、源页图、资源及阶段覆盖保存在JSON与证据笔记。所有计数显式使用本例 $K=128,N=16$，共享128×128仅作方法算例。

% Generated from data/inputs.json raw_evidence; original and adopted values separated.
\begingroup\footnotesize\setlength{\tabcolsep}{3pt}\renewcommand{\arraystretch}{1.08}
\begin{longtable}{@{}>{\raggedright\arraybackslash}p{27mm}>{\raggedright\arraybackslash}p{62mm}>{\raggedright\arraybackslash}p{65mm}@{}}
\caption{参数证据与采用方法。页码均为本地PDF页序；原始值与本文预算分列。}\label{02_sram_dcim:tab:dcim-evidence}\\
\toprule 证据与定位 & 原值、单位和条件 & 采用值与换算理由\\\midrule\endfirsthead
\toprule 证据与定位 & 原值、单位和条件 & 采用值与换算理由\\\midrule\endhead
E01 \texttt{SDCIM-01}\newline p.1 Fig.1；p.2首段 & 物理阵列与完整INT8求值：128×128 bit；128×16 INT8；16项×16输出/拍；64 clock/VMM。28 nm，6T binary SRAM，8-bit signed/unsigned full-precision；输出23 bit & 单原生128×128物理bit，对应K128N16 INT8、64拍完整向量。128物理bit列/8bit权重=16逻辑输出；不聚合闲置容量分片\\\addlinespace[2pt]
E02 \texttt{SDCIM-01}\newline p.2 II-D；p.3 Fig.5 & 完整MAC槽与静态读：16 HCA +16 BFA；8次input-bit展开后换行组；23-bit RCA。静态dual-wordline access；每LBL 16 cell；BFA不在8输入位之间重切WL & 每轮含本地读保持、1-bit输入×8-bit权重、16项归约、符号/移位/累加；无独立read\_ns。完整MAC覆盖静态读与计算；8次行组选择支撑64个bit-round，不新增锁存或重复感测\\\addlinespace[2pt]
E03 \texttt{SDCIM-01}\newline p.4 Fig.10(a)；p.3 III & D6CIM计算时钟：0.9 V: 233 MHz；1.1 V: 正文360 MHz、图363 MHz。28 nm，VMM shmoo；温度未在该图明确报告；不采用最低0.6 V/30 MHz情景 & T\_C=4.3/5/10 ns；4.3 ns由1000/233向上取0.1 ns；5/10由共同数字节拍选择。0.9 V与共同近标称点相邻；不从1.1 V高速点外推0.9 V；三值为预算不是PVT分位数\\\addlinespace[2pt]
E04 \texttt{SDCIM-01}\newline p.1 Fig.1 & 普通SRAM写接口：Write Data[127:0]；Addr[6:0]；CEB/WEB。memory R/W column/row periphery，与VMM控制分开；未独立报告memory write ns & 128 bit=16个完整INT8权重/事务；单分片单行，单更新域。每权重8 binary cell；bitcell互补内部节点不产生额外逻辑payload；不把CIM频率当写频率\\\addlinespace[2pt]
E05 \texttt{CMOS-07}\newline p.2 Fig.2；p.9 Tables II/III & 普通SRAM完整同步周期锚点：455 MHz；48 bit×1k word；2RW 8T；1.05 V typical；25 °C读写功耗测试。28-nm eFlash工艺，高阈值；共同CLK上升沿捕获命令/地址/数据；图示连续write/write/read；频率非单独write-min测量 & 1000/455=2.197802… ns，取2.2 ns为写周期参考下端；5/10 ns为公共节拍预算。以完整时钟和同步写组织支持普通写，不用read access。跨48→128 bit与8T→6T是明确参考移植，需128路写驱动，非同芯片实测配对\\\addlinespace[2pt]
E06 \texttt{CMOS-07}\newline p.9 Table II & read access排除项：1.54 ns。28-nm 48-kbit 2RW SRAM，typical read access & 不用于t\_write或独立DCIM读附加项。read access只覆盖读访问，不是周期，更不是写完成时间\\\addlinespace[2pt]
E07 \texttt{CMOS-03}\newline p.1 II；p.3 IV-C & 6T写成功终点与周期交叉证据：50 FO4$\approx$1 ns；WL=25 FO4；15 fF/128 cell；SA offset0.1 V。HD 28 nm 6T，TT，nominal1 V；瞬态模型；周期末检查内部节点达到目标，支持back-to-back；目标BER<1e-9对应1MB/90\% yield & 仅支持同步完整写及数ns量级；不把1 ns当128-bit D6CIM实测写周期。约半周期WL脉冲$\ne$完整更新；需含节点翻转、保持/关WL和下一周期可读\\\addlinespace[2pt]
E08 \texttt{SDCIM-03}\newline p.1；p.2 Figs.11.7.2-5 & 其他28 nm数字组织：3 ns/clock@0.8 V；8 clock/8-bit乘法；$\Sigma$4/$\Sigma$9/$\Sigma$32 post-sum。dynamic logic；memory和CIM模式分离；lossless bitwise computation；不同post-sum组织 & 仅交叉支持数ns计算时钟和逐位完整精度；不把24 ns直接作为128×128 MVM。3 ns是本宏时钟；输出数、post-sum及本地接口与D6不同\\\addlinespace[2pt]
E09 共享基线\newline JSON共同条件/R0 & 共同逻辑/外围/边界：R1–R8统一逻辑字节、INT8与资源政策；原生K/N显式参数；TD=2/5/10ns、128bit外部编码接口，非固定宏容量。28 nm，0.9 V/25 °C参考，非固定PVT验证；单独立服务单元/更新域；无等面积约束 & 原生K128N16，16项归约，23bit输出；全向量捕获/提交2TD。shared logical参数显式传入；128×128仅公共算例，非硬件约束\\\addlinespace[2pt]
E10 共享基线\newline 通用方法C及R4 & 完整写周期与控制覆盖：T\_front=[2+n\_b-$\chi$]T\_D；n\_b=$\chi$=1；完整周期已含则替代。普通同步memory端口，无另外定义的任务握手，边界在其命令/数据捕获到下一可计算周期 & 主情景Δ\_R=T\_M=2.2/5/10 ns；T\_front额外项为0；外层握手敏感性单列加2T\_D。完整memory周期已含命令解码、数据建立、写驱动和可复用终点；不会在同一阶段上再加T\_W\\\addlinespace[2pt]
\bottomrule\end{longtable}\endgroup

